\chapter{Como as coisas são}

\emph{(Palestra de Dhamma dada a 6\textsuperscript{de} março de 2021)}

Hoje é o dia de observância, a noite da lua minguante. Temos esta oportunidade
esta tarde para refletir sobre o dhamma, como as coisas são. Portanto, para cada
um de nós, a forma como as coisas são neste momento será diferente, não é
verdade? Cada um tem o seu próprio estado de espírito, memórias, pensamentos,
expectativas ou o que quer que seja.

Assim, quando tentamos comparar uma pessoa com outra, ficamos confusos porque
somos todos diferentes. Ao nível dos \emph{saṅkhārās}, dos fenómenos condicionados,
tudo é diferente. Nada pode estabilizar-se numa qualidade ou condição
permanente.

Está além da capacidade. Os \emph{saṅkhārās}, na sua própria natureza, são mudança. E
assim o Buda ensinou: \emph{sabbe saṅkhārā aniccā}, todas as condições são
impermanentes. É assim que as coisas são.

Os \emph{saṅkhārās} são assim. Podem ter qualquer qualidade, alta, baixa, boa ou má,
certa ou errada, seja o que for. Materiais ou mentais, emocionais, eles mudam. A
sua natureza é: anicca -- impermanência, dukkha -- insatisfatória, anatta --
não-eu. Assim, o Buda expôs este ensinamento de forma muito clara. É muito
simples, mas se refletirmos sobre isso, não é um ensinamento que se compreenda
facilmente. Não é uma doutrina budista em que se deva acreditar, porque a crença
apenas nos impede de refletir, assumindo que, como o Buda disse isso, então é
verdade, em vez de pegarmos no que o Buda disse e o usarmos, orientando-nos e
olhando na direção que ele aponta.

Portanto, neste retiro de inverno, este é o último mês do retiro. Neste momento,
neste lugar, é assim. Isto pode parecer muito prosaico e aborrecido, é assim,
mas trata-se de usar palavras para abrir a mente à forma como as coisas são, em
vez de tentar determinar se a forma como as coisas são agora é certa ou errada,
boa ou má, verdadeira ou falsa. E aquilo a que chamamos meditação é, na verdade,
atenção plena.

Existem tantas técnicas de meditação disponíveis agora na Internet e vários
outros professores ensinam estilos diferentes. Mas o ponto principal, porque é
assim que o ensino funciona, depende das palavras. Portanto, não há dois
professores que soem exatamente da mesma forma. A forma como me ouvem, como
recebem o que estou a dizer, será exatamente a mesma de uma pessoa para outra?
Duvido.

Mas a reflexão não é sobre se o que estou a dizer está certo ou errado, não se
trata de tentares acreditar que o que estou a dizer está certo ou de tentares
provar que estou errado, é um convite, um incentivo para refletir sobre a
resposta, a reação que estás a ter neste momento é assim. Seja qual for essa
reação, estado emocional, humor que surge no presente para os indivíduos -- se é
tudo diferente, tudo errado, tudo certo, ou que mistura --, não é essa a
questão; é a forma como as coisas são. Por isso, quando falo sobre o Dhamma,
refiro-me à realidade tal como ela é. Na religião, a religião é a espécie de
superfície exterior de tudo. É a casca da fruta.

Não é o sumo, não é a polpa, mas a superfície externa. Assim, as religiões podem
ser diferentes, têm diferenças, tal como existem muitos tipos de fruta
disponíveis. Alguns preferem uma em detrimento de outra, mas isso não torna uma
melhor do que a outra, é apenas assim que as coisas são.

Ninguém vai exigir que sintamos o mesmo e encontremos uma fruta deliciosa com a
qual todos concordemos. Isso seria tirania. É isso que é a tirania: tens de
acreditar no que eu acredito, no que eu digo. Grande parte do ensinamento
religioso é bloqueada por doutrinas, coisas em que tens de acreditar para seres
um membro funcional de um determinado grupo. Ora, a crença consiste em
compreender conceitos pelos quais nos sentimos atraídos, nos quais estamos
interessados, ou talvez não estejamos interessados neles, mas dizem-nos que, se
não acreditarmos desta determinada forma, então somos pecadores, temos algo de
errado, somos apóstatas e temos de abandonar o grupo. Mas isso não reflete a
realidade, é apenas uma forma de tirania em que uma pessoa determina aquilo em
que todos têm de acreditar, ou um grupo, sem questionar.

Essa é, portanto, uma das razões pelas quais o cristianismo se dividiu em tantas
seitas diferentes, tantos tipos diferentes de grupos cristãos, porque as pessoas
tinham visões diferentes sobre os ensinamentos básicos, a história e assim por
diante da religião cristã. O budismo também pode ficar preso a acreditar apenas
no budismo tibetano ou no budismo zen, no budismo coreano, no budismo tailandês,
no budismo Theravāda. Todos nós podemos sentir que a forma específica que
escolhemos é melhor do que as outras. Refleti sobre isso: este sentimento de que
o que eu tenho é melhor do que o que tu tens é um sentimento de palavras que
surgem e cessam. E essa é a própria natureza de toda a linguagem, de todas as
palavras, por mais nobres, por mais belas que possam ser, ou por mais
mesquinhas, vulgares ou desagradáveis que possam ser.

A única coisa que todas as palavras têm em comum são os seus \emph{saṅkhārās}, as suas
condições, os seus fenómenos. E é por isso que vos encorajo sempre a ouvir as
palavras que vos vêm à mente, sem as julgar. Porque, à medida que aprendeis,
basta sentar-vos quietos num lugar tranquilo por algum tempo, e então vários
pensamentos, memórias e conceitos surgem e desaparecem. Não estão a tentar
controlá-los, como se estivessem a raciocinar logicamente sobre tudo e a julgar
se é certo ou errado, ou verdadeiro ou falso, mas é assim. Portanto, isto é
refletir sobre a natureza das palavras, dos pensamentos, à medida que surgem e
desaparecem, em vez de tentarem descobrir que tipo de pensamentos devem ter e
que tipo de pensamentos são errados e maus, que não deveriam surgir. Se tiveres
pensamentos maus, facilmente assumirás que há alguma fonte maligna em ti ou
algum demónio a tentar tentar-te, ou que és uma pessoa má, porque uma pessoa
boa não teria pensamentos maus.

Então, como vê, isso é proliferação conceptual, em que a mente dá voltas e
voltas sobre o certo e o errado, o bom e o mau. E, claro, presumo que todos aqui
querem ser bons, sabem, em toda a forma monástica. Em todas as formas
religiosas, a ambição de ser uma boa pessoa é um laço comum que partilhamos. Mas
podemos sempre ter bons pensamentos; viver numa comunidade, sabem, tem uma
estrutura com a qual todos concordamos em viver, em render-nos para viver
dentro da estrutura do Vinaya, dos preceitos.

Isso é um acordo para nos juntarmos à comunidade de monges e monjas. Mas isso
diz respeito à ação e à fala, sabem, por isso todos os preceitos do Vinaya dizem
respeito à ação correta e à fala correta, não ao pensamento correto, porque o
pensamento pode ser errado, pode ser mau, pode ser maléfico, assim como bom e o
melhor, o pensamento mais elevado possível que se pode ter em qualquer momento.
Não se consegue sustentá-lo. Pensamentos maus, da mesma forma, não se conseguem
sustentar. Se os observares com um, eles surgem e cessam. Quando lhes resistes,
quando tentas suprimi-los, então proliferas em torno deles e culpas alguém,
quer culpes fontes externas, quer alguma espécie de força maligna no universo
que te está a tentar, é uma forma de expressar isso, ou penses que alguém na
Sangha está a tentar influenciar-te de forma negativa, ou que és simplesmente
uma pessoa má. Mas seja qual for a sua interpretação em relação aos seus
pensamentos maus ou malignos, continua a ser o uso de palavras, a proliferação
de conceitos, seja na língua em que pensa ou na língua que surge na sua
consciência.

Mas quando nos refugiamos no Buda, no Dhamma e na Sangha, Buda, Buddho significa
consciência desperta; não se trata de tentar tornar-se um Buda ou agir como um
Buda, por isso não é um espetáculo em que tentamos vestir uma fantasia para
parecer Budas, rapando a cabeça e vestindo estas vestes coloridas. Não estamos a
tentar tornar-nos Budas ou agir como Budas, mas são os sinais externos, a
casca, a superfície, sobre os quais podemos refletir, como se a túnica fosse
sempre aquela que te lembra que já não és um leigo, porque muitos dos nossos
pensamentos, muitas das nossas reações emocionais, experiências que temos ainda
vêm da época em que éramos leigos, mesmo da época em que éramos crianças
pequenas, adolescentes, jovens adultos e assim por diante. Mas o que é, ou quem
é, que está consciente dos pensamentos como pensamentos, e não o crítico? O
crítico não é quem tu és; a tua posição na vida não é ficar preso numa mente
crítica, que está sempre presa nisto, em ver tudo, o que está certo ou errado em
cada condição, em cada situação, porque todas as condições, todos os \emph{saṅkhārās}
estão a mudar, não se podem sustentar, não se pode sustentar os \emph{saṅkhārās}. Tal
como muitos de nós tivemos insights profundos através de várias formas de
meditação e depois lembramo-nos disso, lembramo-nos que a meditação de ontem
foi realmente boa. Por exemplo, acho que ontem me sentei no templo e tive isto,
o insight mais profundo sobre o dharma.

O que é isso? É uma memória, não é, que surge no momento presente; é assim que
as memórias são impermanentes. No dia seguinte, voltas e sentas-te no mesmo
lugar no templo e fazes as mesmas coisas que te lembras de ter feito no dia
anterior, quando tiveste essa percepção profunda, e o que acontece? A tua mente
fica toda inquieta, negativa, estás a lutar contra isso, queres levantar-te e
ir embora, sentes-te desapontado, esperas que essa memória do dia anterior, a
bem-aventurança que experimentaste no dia anterior, possa perdurar por toda a
tua vida para sempre, esperamos, gostaríamos disso. Sabes que gostaríamos de
viver num estado daquilo a que chamamos bem-aventurança para sempre e sempre,
porque não gostamos de sofrer.

Mas a própria natureza dos \emph{saṅkhārās} é dukkha, ou seja, sofrimento e
insatisfação, porque é assim que os \emph{saṅkhārās} são: não são satisfatórios. Por
mais bons, belos, corretos ou perfeitos que possam ser, acabarão por mudar.
Porque os \emph{saṅkhārās} são impermanentes e insatisfatórios, e é assim que as coisas
são. De quem é a culpa disso, o que achas? É culpa de Deus? Por que é que Deus
não criou \emph{saṅkhārās} permanentemente felizes para nós, para que, uma vez que
tenhamos a percepção, permaneçamos nesse estado de felicidade para sempre, para
além da morte? Porque podemos imaginar, sabes, que a felicidade é um estado
permanente, mas a imaginação que criamos sobre a felicidade continua a ser
palavras, continuam a ser conceitos que crias no momento presente.

Portanto, tentar recordar insights anteriores -- é sofrimento. Quando eu era um
sammanera, há 55 anos, antes de conhecer Luang Por Chah, passei o primeiro ano
como sammanera num mosteiro em Nongkai, que fica em, no nordeste da Tailândia;
é o local por onde se passa para atravessar para o Laos e ir para Vientiane. Fui
ordenado em Wat Sisakit, um dos principais templos no centro da cidade, e o
monge superior enviou-me para um mosteiro de meditação fora da cidade, onde
passei um ano a meditar, tentando, inicialmente, usar métodos que, sabe, não
conseguia, que bloqueavam a minha\ldots{} bloqueavam-me, não conseguia
ultrapassá-los. Eu sentia, sentia que, se não estivesse a usar esse método, não
estava a meditar, sentia-me culpado. Tentei forçar-me a fazer esta, esta
técnica de meditação o dia todo, e quando se está consigo mesmo, sabe, não há
distrações. Não levei nenhum livro, exceto um livro que já dei a muitos de vocês
chamado «A Palavra do Buda», os ensinamentos básicos dos suttas do Tripitaka.

Esse é o único livro que me permiti ler. Assim, nos primeiros três meses, depois
de tentar desesperadamente praticar esta técnica, que me foi ensinada em
Banguecoque antes de me ordenar, e não conseguir, simplesmente desisti, e depois
fiquei sentado ali durante cerca de três meses num kuti sozinho, e elas foram
muito boas para mim, as monjas, as sammaneras que me forneciam comida todos os
dias, tive um bom apoio da comunidade em Nongkai. Portanto, não tinha nada de
que me queixar, não se devia a nada, sabe, de indesejável que estivesse a
acontecer à minha volta no mosteiro. Mas eu tinha 31 ou 32 anos na altura e
tinha passado a vida a reprimir estados negativos.

Por isso, a minha autoimagem era a de que eu era uma pessoa bem-disposta,
basicamente uma pessoa muito bem-disposta, porque sabia como me dar bem na vida
e ser simpático e aberto, e considerava-me um homem adulto bem ajustado. Mas
depois, a viver sozinho sem nada para fazer 24 horas por dia, 7 dias por semana,
exceto praticar esta técnica, não consegui aguentar. E então disseram-me, sabe,
o professor que me ensinou esta técnica disse: «Tenho de continuar a fazer isto
repetidamente até alcançar a iluminação.»

Bem, não estava a funcionar. Então tive esta intuição, uma espécie de intuição,
de que não havia nada que eu pudesse fazer, por isso aprendi a simplesmente
sentar-me e observar. E tanta raiva, ressentimento, medo começaram a surgir na
minha consciência. Olhei para isso como 32 anos de raiva reprimida,
ressentimento, na vida de qualquer pessoa, não é? Muito para se ressentir. A
vida não nos vai tratar de forma justa o tempo todo, nem nos vai valorizar, nem
nada disso. Sabes, a vida tem as suas qualidades de justiça e bondade e também o
seu oposto. Mas disseram-me que a raiva era um pecado. Então, simplesmente,
sabes, quando me sentia zangado, não estava zangado com ninguém no mosteiro,
apenas com ressentimentos de que me lembrava desde a infância até ao tempo em
que era estudante e quando estava no exército. E ao tentar parar esta raiva,
decidi que não ia pará-la, apenas a deixei ir.

E foi o que fiz. E a raiva é um \emph{saṅkhāra}. Sabes, não é permanente. Não é o eu. É
anatta. Então, quando resistia, havia um hábito de uma vida inteira naquela
altura, 30 anos a resistir, a reprimir sentimentos negativos, medo. Sabes, na
minha geração, aos homens americanos dizia-se que não devíamos ter medo de
nada. Os nossos modelos de masculinidade eram cowboys e atletas que pareciam
sempre destemidos, corajosos e que enfrentavam os perigos com grande coragem e
força. Então, é assim que um homem deve ser, de acordo com o condicionamento
cultural.

Mas é assim que as coisas são. Estar sozinho, por minha conta, no escuro, sem
eletricidade, apenas com uma vela ou um candeeiro de parafina, era um verdadeiro
luxo. Num país estrangeiro cuja língua eu não compreendia, sozinho no escuro, na
floresta, como sabem, o medo começou a surgir; eu deixei, permiti que o medo
surgisse, em vez de tentar detê-lo, racionalizá-lo ou justificá-lo. É assim,
o medo é assim, a raiva é assim, e você, de certa forma, abre-se a isso e
simplesmente é o Buddho, a testemunha disso. É o que é no momento. E se não
fizer nada, relaxar e deixar que seja o que é, isso cessará. E é porque essa é a
natureza dos \emph{saṅkhārās}.

Eles têm de cessar. Não existe uma capacidade permanente de um \emph{saṅkhāra} se
sustentar por muito tempo. Então, ao fim de três meses, acordei uma manhã neste
kuti, uma pequena cabana de madeira com telhado de zinco, e todo o lugar estava
simplesmente lindo. Eu estava meio que luminoso, rodeado de luz, e nem
conseguia, tentar ter um pensamento negativo, não conseguia. As memórias
anteriores que costumavam deixar-me irritado e ressentido não faziam efeito,
estavam simplesmente vazias, não conseguiam afetar-me. Já não conseguiam
fazer-me acreditar nelas. Então, pensei que estava iluminado.

Imediatamente pensei que devia ser assim que a iluminação era. E isso durou
cerca de uma semana, e, sabes, eu aproveitei ao máximo. E depois tive de renovar
o meu visto para ficar na Tailândia.

Tive de ir ao posto de imigração em Nongkai para prolongar o visto para ficar na
Tailândia. E foram muito mal-educados comigo nesse posto de imigração. E essa
luminosidade radiante desapareceu completamente. Então, voltei para o mosteiro.

A propósito, eles acabaram por prorrogar o visto. Voltei para o mosteiro e
tentei recuperá-lo. Continuei a tentar recuperar aquele estado de luminosidade,
de êxtase. Tudo era lindo. A floresta e o mosteiro eram lindos. O kuti, o
telhado de zinco e o coração de madeira onde eu vivia eram como um palácio.

Tudo estava banhado de luz e beleza. E era isso que eu queria. Por isso, criei o
desejo de ter essa mesma experiência.

Mas por mais que tentasse recuperá-la, ela não voltava. Então comecei a rever
esta palavra do Buda, as Quatro Nobres Verdades, basicamente o primeiro sermão
do Buda, e comecei a olhar, sabe, a tentar contemplar o sofrimento, a Primeira
Nobre Verdade. E percebi que, sabe, a causa do sofrimento é querer algo que não
se tem. Então comecei a ter a percepção de que não tenho essa felicidade
luminosa e que a quero.

Isso é bhava tanha. É o desejo de obter algo de que me lembro que não tenho
neste momento. Então comecei a despertar para o dhamma, para a forma como as
coisas são. Bhava tanha é assim, querer obter algo de que te lembras ou que
concebes para que possas alcançar a iluminação ou tornar-te um arahant, ou para
te livrares de maus pensamentos, kilesas. Querer livrar-te disso é vibhava
tanha, o desejo de te livrares do que não gostas, do que não queres.

É resistência, repressão. Assim, neste Ensinamento das Quatro Nobres Verdades,
estes três tipos de desejo -- kamma tanha, bhava tanha, vibhava tanha --
despertaram-me verdadeiramente para o Dhamma, apenas ao refletir sobre o
desejo. Porque, na atitude com que cresci na vida familiar cristã americana, o
desejo era visto como algo de mau. Normalmente, significava desejo sexual. Somos
celibatários, não devemos ter desejos sexuais. Ou o desejo no dhamma, pode ser
dividido nestas três categorias, que são muito úteis para perceber quanto da
vossa vida aqui em Amaravati tem a ver com vibhava tanha, o desejo de se livrar
de coisas ou o desejo de obter algo que se quer, de alcançar a iluminação, de se
tornar um arahant. Ora, não há nada de errado em querer tornar-se um arahant ou
alcançar a iluminação.

Já não se trata de bom ou mau, mas é um desejo. E os desejos são impermanentes,
são \emph{saṅkhārās}. Todos os desejos, todo o tanha, são impermanentes e não são o eu.
Assim, estas três características -- anicca, dhukka, annata -- dão-nos esta
informação maravilhosa de que todos os \emph{saṅkhārās}, o que todos partilham, desde o
melhor ao pior, do maior ao menor, é que são impermanentes, insatisfatórios e
não pessoais, não são o eu. Contemple isso: o que quer que pense que é, como se
concebe a si mesmo, isso é um \emph{saṅkhāra}. Por isso, ouça o que pensa que é, o que
acredita que é, seja de forma positiva ou negativa, não importa.

É o Buddho, a consciência, o momento de consciência desperta em que estamos
cientes de que os pensamentos sobre mim, o que penso, os meus sentimentos, o meu
corpo, a minha posição, a minha idade, o meu género, os meus direitos, tudo isto
são pensamentos que surgem e cessam, em vez de conceitos a serem agarrados e
proliferados. Ora, esta é a genialidade do Buda ao dar este ensinamento que é
muito direto. Não é obscuro nem secreto. Portanto, ao transmitir este
ensinamento a monges, monjas e leigos, não se trata de um ensinamento especial e
privilegiado para pessoas muito especialmente ungidas com elevadas qualidades
espirituais.

Está disponível para todos e tem estado disponível há cerca de 2.560 anos no
planeta. Mas grande parte da cultura, da civilização como a chamamos, não se
baseia nesse tipo de sabedoria, não se baseia na sabedoria, baseia-se em
ideais, em como as coisas deveriam ser. Os ideais são, quer consideremos como a
vida deveria ser, como Amravati deveria ser como um ideal\ldots{} como a vida
monástica deveria ser, todos nós podemos imaginar como deveria ser: deveria ser
completamente justa, igualitária, correta, totalmente certa. E se aderirmos ao
culto do budismo Theravāda, então acreditamos que somos melhores do que outras
formas de budismo e outras religiões, porque acreditamos em conceitos sobre o
melhor ou o ensinamento budista, o ensinamento puro, o ensinamento original.

Podemos ser críticos em relação a outras formas de budismo porque ensinam de
maneiras diferentes. Mas tudo isso é proliferação conceptual. É usar a mente
pensante para fazer julgamentos, julgamentos de valor, sobre nós próprios, sobre
as condições em que vivemos, sobre o estado do sistema político e do sistema
social. Assim, todas as guerras, todos os conflitos de que ouvimos falar nos
meios de comunicação social têm a ver com proliferação conceptual. Cada lado
pensa que está certo. E o outro, por não concordar, está errado, sabe, temos uma
mente crítica. Na América, somos muito educados com a ideia de que a democracia
é o melhor e que a América representa a democracia pura. Foi isso que me
disseram quando era jovem. Depois descobrimos todas as coisas antidemocráticas
que acontecem. E não devia ser assim. Não deviam ser hipócritas e deviam estar à
altura deste ideal de democracia. Mas a democracia é um ideal. É uma palavra. É
uma palavra bonita. É muito, sabes, muito inspiradora, pode ser inspiradora.
Como o socialismo. Na América, se disseres que és socialista, és considerado
comunista e traidor da democracia. O socialismo tem uma conotação muito
pejorativa nos Estados Unidos, pelo menos tinha quando eu lá vivia. Mas noutros
países, chamam-se a si próprios sistema socialista, socialistas democráticos e
assim por diante; pode-se chamar ao sistema socialismo democrático, socialismo,
comunismo.

São todas palavras. Mas nenhuma delas está à altura do ideal porque não podem
estar. Porque a vida é assim. As pessoas são como são. Não podemos, não somos
todos Arahants ou bodhisattvas perfeitamente iluminados. O que nós, e, sentimos
é assim, o que nem sempre é certo ou bom, mas é como as coisas são. E esta forma
de refletir ajuda-nos então a aceitar a vida tal como ela flui através de nós.
Sabes, estamos a desenvolver sabedoria.

Estamos a usar a sabedoria com consciência plena. Ficamos desiludidos
com\ldots{} Lembro-me de me juntar a movimentos pela paz quando estava em
Berkeley, no início dos anos 60, porque idolatrava a paz. E eles tinham um
movimento pela paz muito ativo. Havia várias organizações diferentes de
movimentos pela paz em Berkeley na altura, em Berkeley, Califórnia. E a Comissão
de Energia Atómica tinha um escritório em Berkeley naquela altura. Então, íamos
protestar, levando a ciência e dizendo «paz».

E porque considerávamos a Comissão de Energia Atómica, sabe, não pacífica e
exigíamos que se tornassem pacíficos. E enquanto carregava este cartaz a dizer
«paz», olhei para mim mesmo e comecei a pensar: «Não sei do que estou a falar.
Na verdade, não sei o que é a paz. Tenho uma ideia do que é a paz. Mas,
pessoalmente, não sou pacífico como pessoa.»

Eu sou todo, a maior parte do meu estado mental parece não pacífico. E então
pude ver dentro de cada movimento pela paz -- eu participei em dois movimentos
pela paz naquela altura -- e havia muita inveja e conflitos dentro do grupo.
Todos de mente elevada, idealistas, pacifistas, pessoas que exigem paz dos
governos, das instituições políticas, da religião, dos seus companheiros, dos
seus parceiros, dos maridos, das esposas, todos querendo paz. Mas o que é a paz?
E isto, claro, é que, quando compreendemos o sofrimento, começamos a perceber a
paz; a nossa verdadeira natureza é pacífica, é basicamente pacífica. O que não
és são as condições de inquietação que surgem e cessam. Mas subjacente a todas
essas condições tensas, sejam elas quais forem, está a verdadeira natureza da
consciência: a paz, a atenção plena, a consciência.

Nisto podes confiar. Este é o teu refúgio. Não podes refugiar-te nos teus bens
pessoais, nas tuas preferências pessoais, porque eles vão mudar e então haverá
conflito, porque outras pessoas têm apegos diferentes a ideias diferentes.

Por isso, tentar encontrar um conceito, um \emph{saṅkhāra}, em que todos concordemos é
impossível, porque somos todos diferentes nesse nível de \emph{saṅkhārās}. Consegues
realmente mudar a forma como és? Consegues realmente ser outra pessoa? Consegues
simplesmente transformar-te num Arahant iluminado, ou num Buda, ou num
bodhisattva só porque esse é o ideal que defendes? É possível para qualquer um
de nós forçar-nos a tornar-nos perfeitos? É impossível, não é? Porque os
\emph{saṅkhārās} não são perfeitos.

A sua própria natureza é a imperfeição, a mudança. Portanto, os ideais são
ideais. Podes levar o pensamento ao superlativo, ao melhor, à forma mais elevada
possível de pensar, mas isso é um pensamento. Um pensamento é um \emph{saṅkhāra}.

É anicca, dukkha, anata. Tudo muda. Surge e desaparece. Não é possível manter
ideais perfeitos. Podemos apegar-nos a eles e acabar presos a julgamentos sobre
nós próprios, sobre os outros, a julgamentos de valor, porque ninguém consegue
estar à altura do ideal a que nos agarramos. Por isso, muitas vezes ficamos
desiludidos com os sistemas políticos, com a religião, com a meditação. Quantos
de vocês acham que não são bons meditadores? Porque não conseguem atingir o
\emph{samādhi} ou não têm \emph{jhānas}, não são bons meditadores, porque quando se sentam no
templo onde os outros parecem perfeitamente concentrados e em \emph{samādhi}, a vossa
mente vagueia por todo o lado, e pensam que isso é o que vocês são, e pensam que
não são bons meditadores, ou que não conseguem meditar e, ou que não são
realmente budistas, ou criam todo o tipo de pensamentos proliferantes sobre
isso.

O caminho direto, então, não consiste em alcançar \emph{jhānas} e o \emph{samādhi} perfeito,
mas em reconhecer que as condições que estás a vivenciar no presente são o que
são. São como são. Estão a mudar, e a tua relação com elas consiste em
testemunhar a sua mutabilidade, ser paciente, deixá-las ser o que são, o seu
surgimento e cessação. De acordo com a ação e a fala, fazemos o nosso melhor
para nos conformarmos com o Vinaya, com os preceitos. Mas em termos de atividade
mental, hábitos emocionais, não há Vinaya em torno dos hábitos emocionais ou em
torno das memórias, ou pensamentos, ou tendências de caráter, porque todas estas
são condições em mudança, com diferentes qualidades e quantidades em todos nós.

Por isso, reitero quantas vezes já falei assim, mas há apenas um ensinamento
importante. Podes falar sobre todo o tipo de coisas, sobre experiências pessoais
e pontos de vista e opiniões que tens sobre vários outros professores, outras
formas religiosas; posso proferir todas as palavras inspiradoras sobre o budismo
Theravāda, a Tradição da Floresta Tailandesa. Posso inspirar-vos com palavras
inspiradoras, mas são apenas palavras, e por isso esta ênfase na insatisfação
impermanente das palavras não significa que haja algo de errado com elas, mas
sim que são muito limitadas, tal como os \emph{saṅkhārās} não são o que realmente sois,
e quando vos identificais, quando vos agarrais aos \emph{saṅkhārās} que acreditam ser
quem sois, vão ficar infelizes, porque mesmo no seu melhor vão desapontar-vos.
Então, afinal, onde reside a paz? Onde está a paz neste momento? Está aqui no
templo em Amaravati? Podes controlar, sabes, tal como tantos monges que conheci,
e monjas também, que querem condições externas pacíficas. Sabes, sem ruído, sem
conflitos, apenas para viver num estado de bem-aventurança e paz. Ficamos
incomodados porque o ar, uh, o aeródromo de Luton, os aviões a sobrevoar\ldots{}
isso irrita-me\ldots{} a caixa do cão, alguém a cortar a relva, e isso perturba
a minha prática de \emph{mettā} ou a minha paz, porque nós, nós pensamos na paz como
controlar tudo para que nada nos perturbe, nada nos distraia. Bem, não é assim
que as coisas são, a vida, sabes, ter olhos, ouvidos, nariz, língua, corpo,
todas estas são condições inquietas. Não são pacíficas, os olhos não são
pacíficos, o nariz, os odores mudam, a audição, os sons, agradáveis,
desagradáveis, neutros. A experiência sensorial não é pacífica, a sua própria
natureza muda, e é assim que as coisas são. E não querer que seja assim é
torná-la sofrimento; estás a criar sofrimento em torno da forma como as coisas
são, ou, se a tua verdadeira natureza é pacífica, então estás em paz com as
condições mutáveis que estás a vivenciar. Elas são assim. A forma como as coisas
são neste preciso momento, seja o que for que estejas a sentir, seja o que for
que estejas a pensar, é assim. E neste sentido de abertura, abro bem os braços,
como se abraçasse o momento, em vez de tentar controlá-lo para um tipo de
pensamento perfeito, uma equanimidade perfeita, que imagino, ou que me lembro de
ter tido, também me abro. Se há conflito, é assim. Se há caos, é assim.

Se houver ruído, sons cacofónicos, maus odores, coisas desagradáveis de se ver,
é assim. Por isso, isto pode fazer-se em qualquer lugar, quer esteja aqui no
mosteiro no centro de Londres, num engarrafamento, numa reunião com outras
pessoas. Assim, a meditação não se limita apenas a sentar-se no templo em
determinados momentos do dia. Está integrada na forma como nos movemos e
mudamos, no sentar, no ficar de pé, no caminhar, no deitar, na inspiração, na
expiração, no movimento e na mudança dos Sankaras, do Samsara, das condições
mutáveis que todos vivemos através dos nossos corpos, através dos nossos
sentidos. Mas o que é imutável, inamovível, inalterável é um, esta consciência.
Se começarmos a perceber que o nosso refúgio não consiste em agarrar-nos a ele,
mas sim em libertar-nos de tudo. É assim.

Portanto, a verdadeira meditação é o ápice, libertar-se de absolutamente tudo,
o que não é aniquilação. Libertar-se é relaxar, não é tentar alcançar o
\emph{samādhi}, não é tentar obter insight, tentar obter algo que não se tem ou
livrar-se do que se tem e não se quer. Mas é assim, é umas férias relaxantes
para o coração. Luang Po Chah falou-me sobre o que lhe perguntei uma vez há
anos, o que é meditação, \emph{bhāvanā}, se ele poderia definir meditação em tailandês?
E ele disse, em tailandês, «patporn thang Cit tai», o que significa, traduzi
como férias do coração. E pensei, bem, certamente não estou de férias do
coração. Não sei o que isso é. Isto, tentar tornar-me, seguir todas estas
regras do Vinaya e tentar meditar e alcançar o \emph{samādhi} e tentar alcançar os
\emph{jhānas} e tentar alcançar a iluminação é um trabalho árduo. Exige muito esforço.
E, por vezes, simplesmente não consigo fazê-lo.

Mas Luang Po Chah tinha sempre essa sensação de relaxamento aberto de estar no
momento, acontecesse o que acontecesse, porque a sua vida não era só elogios e
flores e adulação, honras e assim por diante. Ele teve de suportar muito stress
e desilusões e condições mutáveis que fazem parte da vida. É assim que as coisas
são naquilo que poderíamos considerar os melhores mosteiros, para não falar de
quaisquer outros.

Portanto, a reflexão desta tarde destina-se apenas a encorajá-los a confiar no
que realmente são, na vossa consciência. Não a tentar tornar-se como outra
pessoa ou como alguma freira ou monge ideal imaginário, ou ser humano iluminado.
E eles podem ser inspiradores, tal como os mestres iluminados são inspiradores
para todos nós. Mas não podem tornar-se um mestre iluminado que compreendem
como um conceito. O mestre é a própria consciência na qual aprendem a confiar
totalmente e a integrar na vossa vida à medida que ela acontece. Sabem, quer se
trate de elogios ou críticas, sucesso, fracasso, felicidade, sofrimento, sucesso
ou fracasso. Porque todas estas são condições mundanas. Estão enumeradas nas
escrituras como os oito Dhammas mundanos. Como boa sorte, má sorte. Todos
queremos boa sorte. Todos queremos sucesso. Todos queremos elogios. Todos
queremos felicidade.

Mas estas palavras, sucesso, boa sorte, felicidade, elogios, são palavras,
palavras positivas que são desejáveis. E não queremos ser menosprezados,
desprezados. Não queremos ser fracassados. Não queremos ser perdedores.

Não queremos ser culpados por coisas. Não queremos ter má saúde. Não queremos
envelhecer. Não queremos adoecer. Não queremos morrer. Temos tantas coisas que
não queremos. Mas estas condições de não querer são \emph{saṅkhārās}. São Dhammas
mundanos. Porque neste momento, este momento só pode ser assim, tal como é.

Por isso, ofereço isto como reflexão. Que todos possam beneficiar com isto e não
se limitem a aceitar, a agarrar-se ao ensinamento, mas que o apliquem à vossa
própria experiência de vida, à medida que vivem momento a momento.
